\documentclass[10pt]{article}

\usepackage[utf8]{inputenc}

\usepackage[T1]{fontenc}

\usepackage{amsmath}

\usepackage{amsfonts}

\usepackage{amssymb}

\usepackage[version=4]{mhchem}

\usepackage{stmaryrd}

\usepackage{mathrsfs}



\begin{document}

Формула (2) справедлива при таких температурах и плотностях, когда среднее расстояние между частицами больше отношения постоянной Планка $h \mathrm{~K}$ модулю средней тепловой скорости



\[

\left(\frac{V}{N}\right)^{1 / 3} \gg \frac{h}{\sqrt{m k T}} .

\]



Частным случаем Б. р. (1) при $U=0$ является Максвелла распределение



\[

f(\boldsymbol{p})=\frac{N}{V}\left(\frac{m}{2 \pi k T}\right)^{3 / 2} \exp \left\{\frac{\boldsymbol{p}^{2}}{2 m k T}\right\}

\]



Функцию распределения (1) иногда называют ра с п р е делением Максвелла - Больцмана, а распре делен ием Больцман а называют функцию распределения (1), проинтегрированную по всем импульсам частиц, представляющую плотность числа частиц в точке $\boldsymbol{r}$ :



\[

n(\boldsymbol{r})=n_{0} \exp \left\{-\frac{U(\boldsymbol{r})}{k T}\right\}

\]



$n_{0}$ - плотность числа частиц, соответствующая точке, где $U=0$. Отношение плотностей числа частиц в различных точках зависит от разности потенциальных энергий, соответствующей этим точкам:



\[

\frac{n_{1}}{n_{2}}=e^{-\frac{\Delta U}{k T}},

\]



где $\Delta U=U\left(r_{1}\right)-U\left(r_{2}\right)$. В частном случае из (4) следует 6 а рометрическая формула, определяющая распределение плотности числа частиц в поле тяжести над земной поверхностью



\[

n(z)=n_{0} e^{-\frac{m g z}{k T}}

\]



где $g$ - ускорение силы тяжести, $m$ - масса частицы, $z-$ высота над земной поверхностью, $n_{0}-$ плотность при $z=0$.



Для смеси газов с различной массой Б. р. показывает, что распределение парџиональных плотностей частиц для каждой из компонент независимо от других компонент. Для газа во вращающемся сосуде $U(\boldsymbol{r})$ есть поле центробежных сил:



\[

U(\boldsymbol{r})=-\frac{m \omega^{2} \boldsymbol{r}^{2}}{2}

\]



где $\omega$ - угловая скорость врашения.



Лит.: [1] Ландау Л. Д., Л ифшиц Е. М., Статистическая физика, ч. 1, 3 изд., М., 1976; [2] Майер Дж., Гепперт-Майе р М., Статистическая механика, пер. с англ., 2 изд., М., 1980. Д. Н. Зубарев.



БОЛЬЦМАНА СООТНОШЕНИЕ - см. Статистический вес. БОЛЬЦМАНА СТАТИСТИКА - статистика систем, содержаших большое число невзаимодействующих частиц (то есть классического идеального газа). Предложена Л. Больцманом (L. Boltzmann, 1868-71). В более общем смысле Б. с.- предельный случай квантовых статистик идеальных газов (Бозе - Эйнштейна статистики и Ферми - Дирака статистики) для газа малой плотности, когда можно пренебречь квантовым вырождением газа, но следует учитывать квантование уровней энергии частиц.



Основа Б. с.- распределение частиц идеального газа по состояниям. Поскольку частицы не взаимодействуют между собой, гамильтониан системы можно представить в виде суммы гамильтонианов отдельных частиц и рассматривать состояния не в фазовом пространстве всех частиц, как в статистич. механике Гиббса, а в фазовом пространстве координат и импульсов одной частицы. Это фазовое пространство разбивается на большое число малых ячеек с таким фазовым объемом $G_{i}$, чтобы каждая из них включала много близких состояний. Это возможно, так как уровни энергии макроскопич. системы расположены очень плотно и стремятся к непрерывному распределению с увеличением числа частиц $N$ и объема тела $V$ (отношение $N / V$ принимается постоянным). Состояние одной частицы соответствует определенной ячейке фазового пространства, а состояние всей системы из $N$ частиц - набору чисел $N_{i}$, характеризующему распределение состояний частиц по ячейкам $G_{i}$. Фазовый объем ячеек выражается в единицах $h^{3}$, где $h-$ постоянная Планка, а число 3 соответствует числу степеней свободы одной частицы. Согласно квантовой механике, координату и импульс частицы можно определить лишь с точностью, допускаемой соотношением неопределенностей, отсюда $h^{3}$ - минимальный размер фазового объема одной частицы (до создания квантовой механики единица фазового объема выбиралась произвольно). Объем $G_{i}$, выраженный в единицах $h^{3}$, имеет смысл максимально возможного числа макроскопич. состояний в ячейке.



В Б.с. предполагается, что частицы распределяются по различным состояниям совершенно независимо друг от друга и что они различимы между собой. Число различных возможных микроскопич. состояний, соответствующих заданному макроскопич. состоянию газа с энергией $\mathscr{\&}$ и числом частиц $\boldsymbol{N}$ (статистический вес $W_{\text {Б }}$ макросостояния по Больцману), определяется числом различных способов, к-рыми можно распределить $N$ частиц по состояниям в ячейках размером $G_{i}$ при $N_{i}$ частиц в каждой ячейке:



\[

W_{\mathrm{Б}}\left\{N_{i}\right\}=N ! \prod_{i} \frac{G_{i}^{N_{i}}}{N_{i} !}, \quad \sum_{i} N_{i}=N, \quad \sum_{i} \varepsilon_{i} N_{i}=\mathscr{E},

\]



где учитывается, что перестановка частиц в пределах каждой ячейки не меняет состояния. При правильном больцмановском подсчете статистич. веса надо, однако, учитывать, что перестановки тождественных частиц не меняют состояния, и поэтому $W_{\text {Б }}$ следует уменьшить в $N$ ! раз:



\[

W\left\{N_{i}\right\}=\prod_{i} \frac{G_{i} N_{i}}{N_{i} !} .

\]



Это правило подсчета состояний лежит в основе Б. с. При таком определении статистич. веса для энтропии системы $S$ получим:



\[

S=k \ln W .

\]



В основе статистич. физики лежит предположение, что все микроскопич. состояния, реализующие данное макроскопич. состояние, равновероятны, поэтому вероятность макроскопич. состояния пропорциональна величине статистич. веса $W$. В статистич. равновесии энтропия максимальна при заданной энергии и числе частиц, что соответствует наиболее вероятному распределению. Его, следовательно, можно найти из условия экстремума $S$ (или $W$ ) при фиксированных $\mathscr{E}$ и $N$. Из этого условия следует Больимана распределение для средних чисел заполнения $i$-го состояния с энергией $\mathscr{~}_{i}$ :



\[

\bar{n}_{i}=\frac{\bar{N}_{i}}{G_{i}}=\exp \left[\frac{\mu-\varepsilon_{i}}{k T}\right]

\]



где $\mu$ - химич. потенциал, $T-$ абсолютная температура. Энтропия идеального газа, подчиняющегося Б. с., равна



\[

S=\sum_{i}\left(N_{i} \ln G_{i}-\ln N_{i} !\right)=-\sum_{i} N_{i} \ln \frac{N_{i}}{G_{i} e},

\]



так как $\ln N_{i} ! \approx N_{i} \ln \left(N_{i} / e\right)$.



Б. с. применима к разреженным атомным и молекулярным газам и плазме, но для плотных газов и плазмы, когда существенно взаимодействие между частицами, надо применять не Б. с., а статистику Гиббса, то есть Гиббса распределение. Б. с. применима к электронам в невырожденных полупроводниках, для металлов надо учитывать вырождение и применять статистику Ферми - Дирака.



Лит.: [1] Л ан д а у Л. Д., Л и ф ш и ц Е. М., Статистическая физика, ч. 1, 3 изд., М., 1976, § 37, 38; [2] Май ер Дж., Гепперт М а й е р М., Статистическая механика, пер. с англ., 2 изд., М., 1980 , гл. 7; [3] 3 о м м е р фе л ьд А., Термодинамика и статистическая физика, пер. с нем., М., 1955, §29. БОЛЬЦМАНА Н-ТЕОРЕМА - одНО из слеДствиЙ Больцмана кинетического уравнения, согласно которому функция



\[

H(t)=\int h(t, \vec{p}, \vec{q}) d^{3} \vec{p} d^{3} \vec{q}=-\int f \ln f d^{3} \vec{p} d^{3} \vec{q}

\]



(где $f(t, \vec{p}, \vec{q})$ - одночастичная функция распределения, удовлетворяющая уравнению Больцмана) является неубы-





\end{document}